% Generated by code/08_generate_results_tables.R
\begin{table}[H]
\centering
\scriptsize
\caption{Reconciliation correction compared with public ETR-parity proxy}
\label{tab:app_profit_correction_benchmark_comparison}
\begin{adjustbox}{max width=\textwidth,center}
\begin{tabular}{@{}lrrrrrr@{}}
\toprule
Parent & Recon. profit (US\$ bn) & Recon. (\%) & Public proxy (\%) & Gap (pp) & Median proxy (\%) & Min. foreign parents \\
\midrule
United States (USA) & 9,286.2 & 54.0 &  3.3 & -50.7 &  0.0 & 19 \\
China (CHN) & 6,174.1 & 32.4 & 13.2 & -19.2 & 11.1 & 18 \\
Japan (JPN) & 2,541.3 & 32.4 & 10.8 & -21.6 & 10.4 & 17 \\
Saudi Arabia (SAU) & 1,342.7 & 32.4 &  0.0 & -32.4 &  0.0 & 15 \\
Canada (CAN) & 1,280.9 & 32.4 & 46.7 &  14.3 & 50.2 & 14 \\
Germany (DEU) & 1,033.1 & 32.4 & 38.6 &   6.2 & 40.3 & 17 \\
Korea, Rep. (KOR) &   836.6 & 32.4 &  2.6 & -29.8 &  0.0 & 16 \\
France (FRA) &   660.7 & 32.4 & 16.0 & -16.4 & 17.3 & 18 \\
Brazil (BRA) &   599.8 & 32.4 & 45.3 &  12.9 & 53.9 & 18 \\
Australia (AUS) &   590.0 & 32.4 &  5.7 & -26.7 &  0.0 & 13 \\
\bottomrule
\end{tabular}
\end{adjustbox}

\vspace{0.5em}
\begin{minipage}{0.96\textwidth}
\footnotesize\raggedright \textit{Notes:} Rows are the ten largest domestic parent cells exposed to the reconciliation correction because observed adjusted profit is unavailable. The public proxy is capped at 55.5\% and uses only country-years with positive domestic profit, non-negative tax accrued, positive foreign-affiliate profit and at least three foreign parent jurisdictions. It is reported as a support diagnostic and is not used as a full-panel point estimator. Source file: data/processed/diagnostics/step12\_profit\_correction\_benchmark\_comparison.csv.
\end{minipage}
\end{table}

